\documentclass[11pt]{article}
\usepackage[a4paper,margin=25mm]{geometry}
\usepackage{amsmath,amssymb,amsthm,mathtools}
\usepackage[T1]{fontenc}
\usepackage[colorlinks=true,linkcolor=blue,urlcolor=blue]{hyperref}
\setlength{\emergencystretch}{2em}
\Urlmuskip=0mu plus 1mu
\newtheorem{theorem}{Theorem}
\newtheorem{lemma}{Lemma}
\newtheorem{proposition}{Proposition}
\newcommand{\F}{\mathbb Q(\sqrt{-3})}
\newcommand{\Bstar}{\beta_*}
\newcommand{\pos}[1]{(#1)_+}
\title{A small bootstrap below the $7/8$ quasi-Riemann bound}
\author{Research note prepared with OpenAI Codex}
\date{9 October 2026}
\begin{document}
\maketitle

\begin{abstract}
We extend the permitted parameter range in the fourth-moment lemma of the
OpenAI cubic-theta argument. The printed proof permits
$13/18\le\kappa\le1$, rather than only $3/4\le\kappa\le1$.
Combining this observation with a rational reparametrization of the original
proof, followed by a second finite bootstrap, yields the zero-free boundary
$43747851/50000000=0.87495702<7/8$.
All new exponent inequalities are checked with exact rational arithmetic.
The argument imports the analytic lemmas of the original paper; this note
is not an independent reproof of those lemmas. The extended energy
induction is verified in Lean; the complete improved zero-free theorem
is not. Its small improvement should be distinguished from the substantial,
still unproved mean-value improvements required for a boundary such as $6/7$.
\end{abstract}

\section{Statement and dependencies}

Let $F=\F$, and let $\Bstar$ be the supremum of $1/2$ and the real parts
of the nontrivial zeros of all primitive finite-order Hecke $L$-functions
over $F$, with arbitrary conductor and height. Principal poles are excluded.
We use the arithmetic conventions, original zero extensions, fixed excluded
prime sets, physical row family, and smooth weights of the OpenAI paper
\cite{openai}. In particular, no coprimality mask is silently deleted.

\begin{theorem}\label{thm:main}
With the analytic inputs proved in \cite{openai},
\[
 \Bstar\le B_2:=\frac{43747851}{50000000}=0.87495702.
\]
Consequently every finite-order Hecke $L$-function over $F$ and every
Dirichlet $L$-function, including the Riemann zeta function, is zero-free
in $\Re s>B_2$. The principal pole at $s=1$ is allowed.
\end{theorem}

The new analytic point is Lemma~\ref{lem:extension}. The other changes are
extensions of parameter-dependent estimates already proved in the source,
and exact finite inequalities. We import the following source results,
identified by their stable \LaTeX{} labels:
\begin{itemize}
\item the $7/8$ theorem and the continuation criterion
\texttt{lem:continuation-criterion};
\item the arithmetic identities, theta reflection, reflected energy
\texttt{lem:reflected-energy}, and additive Gram bound
\texttt{prop:probe-gram};
\item the zero detector \texttt{prop:detector-witness}, marked inverse
moment \texttt{lem:marked}, and sixth-power amplification
\texttt{lem:inverse-amplification};
\item the proof of \texttt{lem:plain}, with the range extension established
below, and logarithmic control \texttt{lem:logarithmic-control};
\item the exact high representation, dynamic local errors
\texttt{prop:probe-errors}, shared contour lemmas, and late height closure;
\item the transfer \texttt{prop:hecke-dirichlet-transfer}.
\end{itemize}
The local source inspected for this note is
\texttt{formalize\_11\_12/papers/paper78.tex}. All references to source
line numbers below refer to that file. Ordinary theorem dependencies on
these inputs are retained. No assertion is made that checking the finite
certificates independently verifies the entire original paper.

\section{Extending the fourth-moment parameter}

Recall the source's fourth-moment estimate. Write $M$ for the effective row
width, $n_1,n_2$ for the two plain-polynomial lengths, $z$ for the total
selected prime length, and $A=n_1+n_2+z$. The source proves
\[
 \sum_{k\in\mathcal R_z}
 |S_{\psi_k}(n_1;W_1)S_{\psi_k}(n_2;W_2)Q_{\psi_k}|^2
 \ll Z^{M+\epsilon}
\]
when $z=0$, without a plain-length restriction, or when
\begin{equation}\label{eq:affine}
 z>0,\qquad A+(6\kappa-1)z\le M.
\end{equation}
For $\kappa<1$ its hypothesis is $\Bstar\le(1+\kappa)/2$.
The prime supports are disjoint, their coefficients belong to the fixed
finite ray-character group, and all the original row, mask, smoothness,
and height hypotheses of \texttt{lem:plain} apply. Its slot mesh is chosen
before the slot count. We do not change those hypotheses.

\begin{lemma}\label{lem:extension}
The entire fourth-moment statement, including its mesh and uniformity
conclusions, remains valid for $13/18\le\kappa\le1$.
\end{lemma}
\begin{proof}
Audit the proof at source lines 12531--14986. Its prime estimate uses
logarithmic control on $\Re s=(1+\kappa)/2+e$ and gives
$|Q|^2\ll Z^{\kappa z+\epsilon}$ with the stated seminorm and height
factors. Logarithmic control is uniform for its real parameter in
$[1/2,1]$; here $(1+\kappa)/2\ge31/36$. The zero-free hypothesis is
unchanged. Thus this step does not require $\kappa\ge3/4$.

At the terminal width $M\le\rho$, equation \eqref{eq:affine} gives
\[
 z\le\frac{M}{6\kappa}\le\frac{3M}{13},\qquad
 \kappa z\le M/6.
\]
This replaces the printed intermediate bound $z\le2M/9$ at lines
12932--12933. Only the unchanged bound $\kappa z\le M/6$ is used in
the terminal exponent, so its extra loss remains at most $\rho/6$.

For $A>5M/6$, the centered comparison at lines 12947--12979 uses
\[
 (6\kappa-1)z\le M-A<M/6.
\]
Since $6\kappa-1\ge10/3$, it follows directly that $z<M/20$.
This replaces the printed chain $z<M/21<M/20$. Put $L=M/4$ as in
the source. Then $A-z>2L$, and its reflected comparison has
\begin{align*}
 A_{\rm comp}&\le3M/2-A+2z+\xi\le23M/30+\xi,\\
 A_{\rm comp}+(6\kappa-1)z
 &\le5M/2-2A+2z+\xi\le14M/15+\xi.
\end{align*}
Both required comparison margins remain at least $M/15$, so the
original requirement $\xi\le\rho/30$ still suffices.

All remaining parameter uses are unchanged. In the child affine defect,
$6\kappa-1\ge0$ permits replacing a surviving length by its parent
length. Removing a greedy prefix of whole slots costs
\[
 \kappa d_z\le F_{\rm act}/6+\kappa\eta\le F_{\rm act}/6+\eta.
\]
The denominator $6\kappa$ is bounded away from zero. The aggregate
error bound uses only $0\le6\kappa-1\le5$. Thus the fixed induction
depth, strict width drop, finite mask treatment, exceptional-row estimates,
and smooth-seminorm propagation keep the same loss allocations.
These are every occurrence of $\kappa$ in the proof. The uniform assertions
at lines 12898 and 14889 consequently hold on the enlarged compact
interval as well.
\end{proof}

\section{The improved row-count formula}

Fix $13/18\le\kappa\le3/4$ with $\Bstar\le(1+\kappa)/2$, and put
\[
 c=\frac1{6\kappa},\qquad \alpha=\frac56,
 \qquad\frac29\le c\le\frac3{13}.
\]
For a zero class write $\delta=2a-1$ and let
$x=q/\delta\in[0,1/2]$ be the normalized mean prime amplitude.
The floor class $\delta_0=1/50$ is handled separately. The source's
detector has inverse and plain lengths $r,m$ with $r+m\ge t$,
$1\le t\le3/2$, up to arbitrarily small requested losses.
At row width one the exact plain capacity is now
$z_P(m)=c(1-2m)$; the inverse capacity stays $z_M(r)=(1-r)/2$.

\begin{lemma}\label{lem:count}
Suppose the physical available prime length exceeds $7/37$ by a fixed
positive amount. Define
\begin{gather*}
 D_x=3-(1+4c)x,\qquad P_x=(2-4cx)(1-x),\\
 J=(\alpha-\delta)D_x+\delta P_x,\\
 R_*(\delta,x)=1-\delta+
 \frac{(\alpha-\delta)\delta P_x}{2J}.
\end{gather*}
For $1/50<\delta\le\alpha$, the source's selected-row argument gives
$\#\mathcal B\ll U^{R_*+\epsilon}(1+T_1)^{A_\eta}$.
All capacity, binning, and witness losses may be made arbitrarily small
before the target is fixed.
\end{lemma}
\begin{proof}
For short witnesses the two ideal exponents are
\[
 A_I(r)=1-\delta[x+(1-x)r],\qquad
 S_t(r)=1-\delta[2cx+(2-4cx)(t-r)].
\]
They cross at
\[
 r_*(t)=\frac{(2-4cx)t+(2c-1)x}{D_x}.
\]
Here $D_x\ge53/26$, $r_*(3/2)=1$, and
\[
 r_*(1)\ge\frac{3-2c}{5-4c}\ge\frac{23}{37},
 \qquad\frac13\le t-r_*(t)\le\frac12.
\]
Indeed $r_*(1)$ decreases with $x$, and its value at $x=1/2$
increases with $c$; the complementary length increases with $t$,
and at $t=1$ its value $(1-2cx)/D_x$ increases with $x$.

Thus inverse capacity is at most $7/37$, while the second marked-width
margin is at least $2r-1\ge9/37$. Plain capacity is at most
$c/3\le1/13<7/37$. The original selection, whole-slot rounding,
zero-capacity treatment, and height estimates therefore apply.
The short crossing value is
\[
 R_{\rm short}(t)=1-\delta+\frac{\delta P_x}{D_x}(3/2-t).
\]
Long inverse witnesses retain the exponent
$L(t)=1-\delta+(\alpha-\delta)(t-1)$.
Choose $t=1+\delta P_x/(2J)$; it lies in $[1,3/2]$ and balances
the two values, giving $R_*$. Positivity of $J$ follows from
$P_x\ge10/13$, $D_x>P_x$, and $0\le\delta\le\alpha$.
The boundary $\delta=\alpha$ uses the same no-slot amplified estimate
as in the source. Since $\kappa$ is fixed from an already established
bound, no $\Delta/4$ capacity replacement penalty is present.
\end{proof}

\section{Changing the probe geometry}

Set
\begin{equation}\label{eq:geometry}
 M=1-\ell,\qquad l_x=(M-b)/2,\quad l_y=(M+b)/2,
 \qquad h=1-l_x+\ell.
\end{equation}
Use exactly the source's finite compensated probe, with total slot length
$\ell$, and with $X=Z^{l_x}$, $Y=Z^{l_y}$. The slots retain their
original physical supports in every rescaled summand. Its signal exponent
is
\[
 C(s)=s+l_x/2-1+h/6=s-(4+b)/6.
\]

\begin{lemma}\label{lem:low}
If $0<\ell\le1/5$, $b>0$, $l_x>\ell$, $l_y>\ell$,
$M>2\ell$, and $l_y-\ell-11b/6>0$, the normalized probe satisfies
\[
 |J_\eta(Z)|\ll_{\eta,\epsilon}Z^{C(B)+\epsilon},
 \qquad B=11/12-\ell/4.
\]
\end{lemma}
\begin{proof}
For a rescaled subset of total length $d\in[0,\ell]$, put
$M'=M-2d$ and $\ell'=\ell-d$. The reflected-energy calculation in the
proof of \texttt{lem:probe-row-norm} gives
\[
 E_{\rm ref}\le M'+\pos{(1+3\ell'-2M')/4}+o(1)
 =M'+\pos{(5\ell-1+d)/4}+o(1).
\]
Specifically, the source estimate $T_d\le H-3d+o(1)$ uses only
$M+\ell=1$; it therefore still makes the maximum in the reflected
energy equal to $H+o(1)$. Its two remaining branches give the displayed
bound. One must retain this extra energy term: it need not vanish for
every rescaled subset.

The Gram hypothesis makes $P_a^2/Y'\ll P_a^{1/6}$ and hence supplies
the same additive norm as in the source. Cauchy--Schwarz, tuple counting,
and the rescaling coefficient then give an extra exponent
\[
 f(d)=-d+\pos{5\ell-1+d}/8.
\]
Since $\ell\le1/5$, $f(0)=0$, and both of its linear slopes,
$-1$ and $-7/8$, are negative. Thus $f(d)\le0$ everywhere.
The modified probe is bounded by $Z^{l_x/2+b/12+\epsilon}$.
Its principal normalizer has only a subpower reciprocal, as in the source,
and $l_x/2+b/12=C(11/12-\ell/4)$.
\end{proof}

For a target $B\in[437/500,7/8]$, the source's second Euler region may
be extended to $\Re s\ge B$, retaining $\Re w\ge19/20$ and
$\Re z\ge33/200$. Its good-prime defect exponents
$4-6\Re s-6\Re z$ and $1-\Re s-\Re w-6\Re z$ remain strictly
below $-1$. All local denominators remain uniformly nonzero.
The four principal error exponents
\[
 -\Re s,\quad-6\Re z,\quad4-5\Re s-6\Re z,
 \quad1-\Re w-6\Re z
\]
are at most $-B$. Thus the principal slot multiplier is
$-1+O(q_p^{-B})$, and its approximation margin may be any
$0<\mu<B\min_i\ell_i$. The unselected principal Euler correction
is holomorphic and can be made within $1/2$ of one by enlarging the
fixed prime cutoff. The same principal normalization and residues apply.

The shared contour statements formally use $\sigma_0\ge7/8$.
Their proofs apply at $\sigma_0=B$ after this Euler extension.
On the small-row line replace their $\mathcal D_1(3/8)$ by
$\mathcal D_1(\epsilon_0)$, where
$0<\epsilon_0\le B-1/2$ is fixed. No contour crosses a new pole.
The absolute local estimates remain $G_p=O(1)$ off the row and
$G_p\ll q_p^{1/2}$ on its support. The principal residue margins
remain $l_y/20$ and $h/600$.

\section{Exact endpoint and remaining ranges}

The high exponent relative to $C(B)$ is
\begin{align}\label{eq:high}
 E_B(d;\delta,q,R)
 &=a-B+h(17/50-1/6)-a l_y-\ell/2+q\ell
       +d(R+\delta/2-17/50),\\
 &\hspace{32mm}a=(1+\delta)/2.\notag
\end{align}
At $d=h$, $q=x\delta$, $R=R_*$, it becomes
\begin{equation}\label{eq:endpoint}
 E_B(h)=\frac{(1+\delta)(1-l_y)}2-\frac\ell2+x\delta\ell
       +h(R_*+\delta/2-1/6)-B.
\end{equation}

Start from the source theorem $\Bstar\le7/8$. The rational geometry
\[
 B_0=\frac{34999}{40000},\qquad b_0=\frac18,
 \qquad\ell_0=\frac{5003}{30000},\qquad\kappa_0=\frac34
\]
gives $\Bstar\le B_0$ by the preceding estimates and the same final
continuation argument. At this geometry the exact certificate
\texttt{certificate.py} proves $E_{B_0}(h)\le-9/100000$ on the full
rectangle $[0,5/6]\times[0,1/2]$, and verifies all remaining strict
margins. This is also the rational extension discussed in \cite{liu}.
The use of $\kappa_0=3/4$ is justified by the already proved $7/8$
bound, so this first step has no circular hypothesis.

Now use $\Bstar\le B_0$ and set
\begin{gather}\label{eq:stage1}
 \kappa_1=2B_0-1=\frac{14999}{20000},\qquad
 c_1=\frac{10000}{44997},\qquad
 B_1=\frac{17499141}{20000000},\\
 b_1=\frac{6170279}{50000000},\qquad
 \ell_1=\frac{2502577}{15000000},\notag\\
 l_{x,1}=\frac{106463393}{300000000},\quad
 l_{y,1}=\frac{143485067}{300000000},\quad
 h_1=\frac{81196049}{100000000}.\notag
\end{gather}
Here $13/18<\kappa_1<3/4$, so Lemmas~\ref{lem:extension} and
\ref{lem:count} apply. The formula $B_1=11/12-\ell_1/4$ is exact.

\begin{lemma}\label{lem:certificate}
At \eqref{eq:stage1}, on
$0\le\delta\le5/6$, $0\le x\le1/2$,
\[
 E_{B_1}(h_1)\le-10^{-8}.
\]
All hypotheses of Lemma~\ref{lem:low}, the strict supply conditions
$\ell_1/h_1>7/37$ and $5\ell_1-h_1>0$, and the Euler conditions
above hold. The floor and intermediate row ranges have a strict
saving relative to $C(B_1)$. The source's absolute small-row estimate
has its own strict saving relative to $C(\Bstar)$.
\end{lemma}
\begin{proof}
This is a finite rational calculation, implemented in
\nolinkurl{research_audit/bootstrap_certificate.py}. To specify the
nonlinear certificate completely, form the bivariate polynomial
\begin{align*}
 N(\delta,x)={}&J\bigl[B_1-(1+\delta)(1-l_{y,1})/2+\ell_1/2
 -x\delta\ell_1-h_1(5/6-\delta/2)-10^{-8}\bigr]\\
 &-\frac{h_1}{2}(\alpha-\delta)\delta P_x.
\end{align*}
Thus $N=J(-E_{B_1}(h_1)-10^{-8})$ exactly. The checker proves
$J\ge174985/269982>0$. It partitions the rectangle into 24 rational
subrectangles and verifies that every Bernstein coefficient of $N$ on
each is nonnegative. The Bernstein basis is nonnegative and sums to one,
so these checks prove $N\ge0$ everywhere. All subdivision endpoints,
coefficients, and sign decisions use rational arithmetic; no sampled
positivity or floating-point optimization is used for this inference.
The source code can emit the entire list of rectangles and exact lower
coefficients as a separate certificate.

For completeness, the auxiliary checks include
\[
 \frac{\ell_1}{h_1}-\frac7{37}
 =\frac{13344541}{819341949}>0,\qquad
 5\ell_1-h_1=\frac{6669553}{300000000}>0.
\]
The floor uses $\delta_0=1/50$, $q=\delta_0/2$, $R=1$ and has
\[
 -E_{B_1}(h_1)=\frac{2939583}{500000000}>0.
\]
For $d\in[1/100,1/2]$, the source's no-slot choice $t=1$ gives the
common bound $R=76/75-2\delta/3$ including the floor. Its slope in
$d$ is positive. The two endpoint values in $\delta$ at $d=1/2$
satisfy
\[
 -E_{B_1}(1/2;1/50,1/100,R)=\frac{2148926943}{10000000000},
 \qquad
 -E_{B_1}(1/2;5/6,5/12,R)=\frac{239232781}{90000000000}.
\]
Both are positive, and this exponent is affine in $\delta$.
For $d\le1/100$, the absolute outer-row lemma gives, directly relative
to $C(\Bstar)$, the saving
\[
 -\bigl[h_1(17/50-1/6)-l_{y,1}/2+63/5000\bigr]
 =\frac{1287029401}{15000000000}>0.
\]
The script verifies all other listed length, Gram, residue, and Euler
inequalities by direct rational subtraction.
\end{proof}

After proving $\Bstar\le B_1$, perform a second finite bootstrap with
\begin{gather}\label{eq:stage2}
 \kappa_2=2B_1-1=\frac{7499141}{10000000},\qquad
 c_2=\frac{5000000}{22497423},\qquad
 B_2=\frac{43747851}{50000000},\\
 b_2=\frac{24681137}{200000000},\qquad
 \ell_2=\frac{6256447}{37500000},\notag\\
 l_{x,2}=\frac{425853437}{1200000000},\quad
 l_{y,2}=\frac{573940259}{1200000000},\quad
 h_2=\frac{324784289}{400000000}.\notag
\end{gather}
The same exact checker, with these arguments, gives
\begin{equation}\label{eq:stage2-endpoint}
 E_{B_2}(h_2)\le-10^{-9}
 \quad(0\le\delta\le5/6,\ 0\le x\le1/2).
\end{equation}
The polynomial is the one in the proof of Lemma~\ref{lem:certificate},
with subscript $1$ replaced by $2$ and margin $10^{-9}$. Its certificate
has 26 rational subrectangles, and
$J\ge87487115/134984538>0$. The auxiliary inequalities all pass again.
For explicit representative margins,
\[
 \frac{\ell_2}{h_2}-\frac7{37}=\frac{587163179}{36051056079}>0,
 \qquad5\ell_2-h_2=\frac{26678653}{1200000000}>0.
\]
The floor, smaller intermediate, and small-row savings respectively are
\[
 \frac{58789949}{10000000000},\qquad
 \frac{956906941}{360000000000},\qquad
 \frac{5148114961}{60000000000}.
\]
Thus the entire conclusion of
Lemma~\ref{lem:certificate} holds for stage two as well.

\section{Continuation and conclusion}

\begin{proof}[Proof of Theorem~\ref{thm:main}]
First obtain the rational preliminary bound $\Bstar\le B_0$ as described
above. We explain the continuation step at either $j=1$ or $j=2$,
assuming the already proved $\Bstar\le B_{j-1}$. Suppose for
contradiction that $\Delta=\Bstar-B_j>0$.
Use the corresponding geometry \eqref{eq:stage1} or \eqref{eq:stage2}
and its one fixed value $\kappa_j=2B_{j-1}-1$.
The zero-free hypothesis of the extended plain moment is justified by
$\Bstar\le B_{j-1}=(1+\kappa_j)/2$, independently of the contradiction.
The detector's actual classes satisfy $\delta\le2\Bstar-1\le\kappa_j$;
certifying the larger rectangle up to $5/6$ is therefore sufficient.

Lemma~\ref{lem:low} bounds the normalized physical probe by
$Z^{C(B_j)+\epsilon}$. Lemma~\ref{lem:certificate} and its stage-two
version control the central and smaller nonprincipal rows.
For $1/2\le d\le h_j$, the frequency
slope is positive because
\[
 R_*+\delta/2-17/50\ge33/50-\delta/2>0.
\]
Thus the bound at $h_j$ bounds the entire interval. Write $m_1=10^{-8}$,
$m_2=10^{-9}$ and choose
$0<\zeta<\min\{5\ell_j-h_j,m_j/8\}$, decreasing it further if needed.
Extend to $h_j+\zeta$: slopes are bounded by two, so its cost is less
than $m_j/4$, and $h_j+\zeta<5\ell_j$ leaves prime supply greater
than $1/5>7/37$. Rows beyond $Z^{h_j+\zeta}$
are negligible by the source's absolute large-row contour, choosing its
fixed $z$ line sufficiently far right after choosing $\zeta$.

Subtract $C(\Bstar)-C(B_j)=\Delta$ from the floor, intermediate, and
endpoint exponents computed relative to $C(B_j)$. The absolute small-row
lemma already supplies its saving relative to $C(\Bstar)$, so no second
subtraction is made there. There is no capacity replacement cost.
The order of the remaining choices matters.
First choose the moment, capacity, and rounding reserves relative to the
fixed high-side margin and $\Delta$, without referring to a principal
slot approximation margin. Choose their permitted slot mesh, then an even
slot count and disjoint prime windows of total length $\ell_j$.
Lemma~\ref{lem:extension} gives uniformity before this slot count is fixed.
Now choose $0<\mu<B_j\min_i\ell_i$; the positive principal contour
margins are $l_{y,j}/20$ and $h_j/600$. Only after $\mu$ exists choose
the remaining amplitude-bin, contour, normalization, and detector losses
small relative to it and to the earlier margins and $\Delta$. All these
real choices precede the target character. This prevents a circular
choice of the mesh through a margin that depends on the later slot count.

For each fixed target choose its allowed arithmetic data and internal
seminorm orders. The source's late height closure then chooses the
positive height exponent, external tail order, and final threshold in
that order. The height cutoff appears only in estimates, not in the
physical probe or its principal signal. Consequently, for some
$0<\omega<\Delta$ and $\sigma>0$ independent of the target,
\[
 |J_\eta(Z)|\ll_\eta Z^{C(B_j)+\omega},\qquad
 |J_\eta(Z)-f_\eta(Z)|\ll_\eta Z^{C(\Bstar)-\sigma},
\]
where
\[
 f_\eta(Z)=\frac1{2\pi i}\int_{(2)}
 Z^{C(s)}e^{(s-5/6)^2}\frac{H_\eta(s)}{L_F^S(s,\eta)}\,ds.
\]
The extended Euler analysis makes $H_\eta$ holomorphic on
$\Re s>B_j$ and within $1/2$ of one. These are exactly the source
continuation criterion's hypotheses at $B_j$, contradicting
$\Bstar>B_j$. Hence $\Bstar\le B_j$. Apply this finite argument first
with $j=1$ and then with $j=2$ to obtain $\Bstar\le B_2$.

Restoring finite Euler factors does not introduce zeros in this positive
half-plane. For a Dirichlet character $\chi$, the norm-pullback factorization
\[
 L_F(s,\chi\circ N)=L(s,\chi)L(s,\chi\chi_{-3})
\]
then gives the Dirichlet conclusion, with the usual principal poles
allowed, exactly as in the source's transfer proposition.
\end{proof}

\section{What this result does and does not establish}

The gain over $7/8$ is $0.00004298$. The bound is also slightly below
$(1507-2\sqrt{921})/1653=0.874957069799\ldots$, the optimum reported
for the fixed $\kappa\ge3/4$ parameter system in \cite{liu}.
That paper already flags the possibility of a tiny improvement on allowing
$\kappa<3/4$; this note supplies a proof amendment and a conservative exact
rational realization of that observation. Prior work by LiXiang
\cite{lixiang} also describes this feedback and reports the limiting
value $0.874957019420098946\ldots$, slightly better than our conservative
rational bound. Its repository explicitly states that the improved Lean
conclusion still assumes \texttt{ZetaInverse.ArithmeticProbeObligation};
the lower-$\kappa$ energy theorem is formalized, while the complete physical
probe obligations remain open there. We do not import that obligation as
a proved theorem: our new analytic proof uses the concrete source
reflected-energy, Gram, local-error, and contour lemmas with the amendments
given above. No claim of priority is made.

This is an analytic research derivation using the original paper's inputs,
with machine-checked exact arithmetic for its new inequalities. A separate
Lean overlay verifies the full energy induction with the weaker floor
$\kappa\ge13/18$, retaining the hypothesis $2\Bstar-1\le\kappa$.
All 63 affected modules compile with \texttt{autoImplicit=false};
the exported certified-band and terminal-certificate endpoints depend
only on \texttt{propext}, \texttt{Classical.choice}, and \texttt{Quot.sound}.
The endpoint is \nolinkurl{research_audit/kappa_formal/KappaFormal.lean};
the exact theorem signatures and axiom audit are recorded in
\nolinkurl{research_audit/kappa_formal/logs/endpoint-axioms.log}.
This verifies the energy-induction extension, not the full assembly of the
improved moment, physical probe, and zero-free theorem in Lean. Nor does
this note independently validate every upstream imported theorem.
It proves neither the $6/7$ mean-value hypothesis nor GRH. In particular,
a lower bound for one exponent model is not an impossibility theorem for
all modifications of the cubic-theta method or for future research.

There is nevertheless a useful exact obstruction for the unchanged direct
majorant. Write $M=l_x+l_y$, $b=l_y-l_x$, and suppose its Gram and
energy exponents satisfy
\[
 g\ge b/6,\qquad e\ge0,\quad
 e\ge(1+3\ell-2M)/4,\quad e\ge M+\ell-1.
\]
The empty-subset low bound then requires
$B\ge5/6+M/12-\ell/6+e/2$. The identity
\begin{align*}
 \frac M{12}-\frac\ell6+\frac e2-\frac1{30}
 ={}&\frac15\left[e-\frac{1+3\ell-2M}{4}\right]
 +\frac1{60}\left[e-(M+\ell-1)\right]+\frac{17}{60}e
 \ge{}&0
\end{align*}
therefore forces $B\ge13/15$. This concerns the sufficiency of that
particular direct majorant. It already blocks $6/7$ and $1/2$ within
that unchanged majorant, independently of any improvement in the row
counts. It does not exclude changing the majorant, retaining cancellations
discarded by it, or replacing the physical probe.

\subsection*{The exact barrier including parameter feedback}

The preceding $13/15$ bound uses only the direct majorant. Retaining also
the existing row-count formula gives a stronger, precise limitation.
This is a statement about the following sufficient exponent package:
retain the direct majorant just displayed, require the high bound
\eqref{eq:endpoint} for every permitted zero/amplitude class, and use a
count exponent $R\ge R_*(\delta,x;c)$ from Lemma~\ref{lem:count}.
Prime-supply restrictions and positive reserves may be retained, since
they can only weaken the count. The geometry has $h>0$, but need not
satisfy $M+\ell=1$ for this calculation. A genuinely stronger count,
cancellation between physical rows, or independent exclusion of a critical
class is outside this package.

\begin{proposition}\label{prop:feedback-barrier}
Within this exponent package, reparametrization and repeated feedback of
proved zero-free bounds into the same fourth moment cannot certify
$\theta<\theta_*$, where
\[
 \theta_* =\frac{11}{12}-\frac{\lambda_*}{4},\qquad
 657\lambda_*^3-954\lambda_*^2+21\lambda_*+20=0,
 \quad \frac16<\lambda_*<\frac15.
\]
The specified root is unique, and exact rational sign evaluations give
\[
 0.87495701942009894<\theta_*<0.87495701942009895.
\]
The lower bound is sharp for the three-constraint relaxation in the
proof. This proposition does not assert sufficiency at the exact endpoint
for every estimate of the full analytic argument.
\end{proposition}
\begin{proof}
The direct majorant already forces $\theta\ge13/15$, so consider only
$13/15\le\theta\le7/8$. Under a contradiction $\Bstar>\theta$, an
available fourth moment with $\Bstar\le(1+\kappa)/2$ requires
$\kappa>2\theta-1$. Consequently its capacity is at most the optimistic
limiting value
\[
 c\le C(\theta):=\frac1{6(2\theta-1)},\qquad
 \frac29\le C(\theta)\le\frac5{22}.
\]
The case $\kappa=1$ also respects this inequality. It holds at each
class even if different valid moment parameters are selected by class.

At $x=1/2$, put
$D=5/2-2c$, $P=1-c$, $J=\alpha D-\delta(3/2-c)$.
Direct differentiation gives
\[
 \frac{\partial R_*}{\partial c}
 =-\frac{\delta(\alpha-\delta)^2}
 {[\alpha(5-4c)+\delta(-3+2c)]^2}\le0.
\]
Furthermore
\[
 R_*(\delta,1/2;c)-\frac23=\frac{Q_c(\delta)}{36J},\quad
 Q_c(\delta)=(36-18c)\delta^2+(57c-78)\delta+25-20c.
\]
For $2/9\le c\le5/22$,
$Q_c(1/3)=3(1-c)>0$, $Q_c(1/2)=4c-5<0$, and
$Q_c'(\delta)\le39c-42<0$ on $[1/3,1/2]$.
Let $\delta(c)$ be its unique root in that interval and take
$\delta=\delta(C(\theta))$. It is an admissible test class: it exceeds
the detector floor and is below every possible $\kappa$, since
$\kappa\ge2\theta-1\ge11/15>1/2$. The optimistic exponent at this
class is $2/3$; monotonicity implies every permitted capacity has
$R\ge2/3$ there.

The two relaxed low constraints are
\[
 \theta\ge L_0=\frac56+\frac M{12}-\frac\ell6,
 \qquad
 \theta\ge L_K=\frac13+\frac{7M}{12}+\frac\ell3.
\]
Since $h>0$, substituting $x=1/2$, $R=2/3$ in the high bound gives
\[
 \theta\ge H_\delta=1+\delta-\frac{1+\delta}{2}M+\delta\ell.
\]
Its coefficient of $b$ cancels. The positive weights
\[
 w_0=\frac{4+18\delta}{9+18\delta},\qquad
 w_K=\frac2{9+18\delta},\qquad
 w_H=\frac3{9+18\delta}
\]
sum to one and cancel both $M$ and $\ell$. Hence necessarily
\[
 \theta\ge\frac{7+18\delta}{9+18\delta},
 \qquad\delta=\delta(C(\theta)).
\]
The right-hand function is strictly increasing in $\delta$. Define
\[
 d(\theta)=\frac{9\theta-7}{18(1-\theta)},\qquad
 \lambda=\frac{11}{3}-4\theta.
\]
Here $d(\theta)\in[1/3,7/18]$ and $\lambda\in[1/6,1/5]$.
The necessary inequality is equivalent to
$d(\theta)\ge\delta(C(\theta))$, or
$Q_{C(\theta)}(d(\theta))\le0$. The exact substitution identity is
\[
 Q_{C(\theta)}(d(\theta))
 =\frac{-3P(\lambda)}{(5-3\lambda)(1+3\lambda)^2},\qquad
 P(\lambda)=657\lambda^3-954\lambda^2+21\lambda+20.
\]
Its denominator is positive. Thus $P(\lambda)\ge0$ is necessary.
Now $P(1/6)=1/24>0$, $P(1/5)=-1088/125<0$, and throughout the interval
\[
 P'(\lambda)=1971\lambda^2-1908\lambda+21
 \le1971/25-1908/6+21=-5454/25<0.
\]
There is one root $\lambda_*$, and necessity becomes
$\lambda\le\lambda_*$, that is $\theta\ge\theta_*$.
The equivalent cubic is
$7884\theta^3-18819\theta^2+14643\theta-3686=0$.
The derivative identity, dual cancellations, exact substitution multiplier,
and rational isolating signs are independently verified by
\texttt{research\_audit/feedback\_barrier\_certificate.py}, using only
exact polynomial and rational arithmetic.

At the limiting capacity and root, the relaxation has equality at
$\ell=\lambda_*$, $M=1-\lambda_*$:
$L_0=L_K=H_\delta=\theta_*$. This proves the stated sharpness for
the relaxation. The scalar $\lambda$ used to eliminate the target is not
an assumption about the free geometric $\ell$ before this equality case.
\end{proof}

This barrier concerns the uniform sufficient bounds, not the existence
of zeros in the critical class used to test them. It rules out further
gains from geometry and this parameter feedback alone. It does not rule
out stronger mean values, an altered probe, or treatments preserving
cancellations discarded by the current majorants.

\subsection*{The supplied Gram-splitting route}
The prior note \nolinkurl{gram_split/note/theta_split_note.tex} correctly
retracts its stronger unconditional claims. Reflection contributes a
row phase of the form
\[
 \zeta_m=(\text{fixed ray factor})\,
 \mu(m)\overline{\alpha(m)}\,\overline{\gamma_2(m)}
 \quad(m\text{ squarefree}).
\]
Consequently its relevant correlation is a pairing of $\pi=\rho\zeta$
with the reflected row, rather than a pairing of $\rho$ alone. An estimate
for the latter sequence cannot be substituted for an estimate of the
former. The required saving for the correct sequence, including every
marked block if slots are retained, remains unproved in that route.
The new bootstrap above does not use that saving.

The norm information alone cannot repair this step. Abstract vectors
$u=Q^{5/12}v$ and $w=Q^{1/2}v$, with $\|v\|_2=1$, satisfy
$\|u\|_2^2=Q^{5/6}$, $\|w\|_2^2=Q$, and attain
$|\langle u,w\rangle|=Q^{11/12}$. Thus those two bounds alone imply
no power saving. Additional arithmetic cancellation is necessary.

Even after removing all additional nonnegative losses, the residual
Cauchy--Schwarz template still requires
\[
 B\ge\frac56+\frac{l_x-\ell}{6}+\frac{g+e}{2}>\frac56
 \quad\text{if }l_x>\ell,\quad g,e\ge0.
\]
This is a second, weaker but more robust, obstruction to $1/2$ within
that precisely specified residual estimate. It does not forbid treating
the whole residual by a new cancellation argument.

\subsection*{The alternate $11/12$ extraction route}
In the companion paper \cite{alternate}, a mean square over rows of
length $D^\kappa$ combined with prime sixth-power extraction gives the
comparison exponent
\[
 B_{\rm prime}(\kappa)
 =\max\left\{\frac12+\frac{5\kappa}{12},\ 1-\frac\kappa6\right\}.
\]
Even if the required mean square were available at any $\kappa\ge0$,
this extraction has minimum $6/7$, achieved at $\kappa=6/7$.
Indeed the convex combination with weights $2/7$ and $5/7$ of the
two displayed affine functions is identically $6/7$.
The original theorem supplies the mean square only beyond the diagonal,
$\kappa>1$, yielding $11/12$ in the limiting exponent.
Integer sixth-power inversion can avoid the second term, but then
approaching $1/2$ requires proving mean-square bounds for arbitrarily
short row families. Neither the cited proof nor the present bootstrap
establishes those estimates. This is a missing input, not a proof that
such estimates are false or equivalent to GRH with the same uniformity.

Finally, the older analysis describes perfect cubic large-sieve
orthogonality as an open conjecture. That description is obsolete:
\cite{nonorth} proves non-orthogonality for the cubic and quartic families.
Its sextic prediction is a conjecture; it cannot itself be used as a
proved sextic obstruction or as a cancellation theorem for the specific
coefficient sequence $\pi$ above.

\begin{thebibliography}{9}
\bibitem{openai} OpenAI,
\emph{The Quasi-Riemann Hypothesis: A Zero-Free Half-Plane
$\Re s>7/8$}, 30 September 2026.
Local source inspected: \texttt{formalize\_11\_12/papers/paper78.tex}.
\href{https://github.com/openai/math/blob/main/preprints/The-Quasi-Riemann-Hypothesis-September-30-2026/paper.pdf}{Public manuscript}.
Stable source labels are used in this note to distinguish theorem statements
from the extensions of their proofs required here.
\bibitem{liu} Baiying Liu,
\emph{Slightly Improved Zero-Free Half-Planes for the Quasi-Riemann
Hypothesis}, 8 October 2026. Local text:
\texttt{six\_sevenths/lit/liu1.txt} and
\texttt{six\_sevenths/lit/liu2.txt}.
\url{https://arxiv.org/abs/2610.12234}.
\bibitem{lixiang} LiXiang,
\emph{RH-Zero-Free-Formalization}, source repository and README,
\url{https://github.com/lixiang90/RH-Zero-Free-Formalization}.
The repository distinguishes the formalized lower-parameter energy
theorem from the conditional improved-bound theorem.
\bibitem{alternate} OpenAI, \emph{The Quasi-Riemann Hypothesis},
alternate $11/12$ proof, 5 October 2026.
\href{https://github.com/openai/math/blob/main/preprints/The-Quasi-Riemann-Hypothesis-October-5-2026/paper2.pdf}{Public manuscript}.
\bibitem{nonorth} A. de Faveri, A. Dunn, and J. Hoffstein,
\emph{Non-orthogonality of the cubic and quartic large sieves via
Rankin--Selberg}, 2026. \url{https://arxiv.org/abs/2607.07911}.
\end{thebibliography}
\end{document}
